\section{Introduction}
\label{sec:intro}

Urban air mobility (UAM) proposes to move passengers across a metropolitan area in
small electric aircraft flying at low altitude~\cite{cohen2021uam}. The technical
obstacles to that vision are increasingly well understood; the binding constraint on
deployment is turning out to be a social one. Community noise is repeatedly identified
as the factor most likely to determine whether a UAM network is permitted to operate at
all, and the response of a community to a new noise source depends heavily on
non-acoustical factors that a purely acoustic analysis does not
capture~\cite{nasaUAMnoise2020,kim2022indoor}.

The distributional question is unavoidable because the burden is spatial. Aircraft must
fly \emph{somewhere}, the vertiports are fixed, and the communities lying between them
absorb the noise. Which communities those are is a choice made by whoever writes the
routing objective. It is also a choice made against a background that is already
unequal: in the contiguous United States, noise exposure is systematically higher in
census block groups with larger minority and lower-income populations, and disparities
in civil aircraft noise exposure specifically persist across
years~\cite{casey2017noise,simon2022sociodemographic}. A UAM network laid over that
background either reproduces the existing pattern or deliberately does not.

Existing noise-aware routing formulations answer the distributional question implicitly,
and almost always in one of two ways: minimise total or average noise, or hold every
receptor below a threshold. Both are defensible. Neither is neutral, and neither is the
only defensible option. Political philosophy has spent decades articulating what a fair
division of a burden might mean, and it does not converge on a single
answer~\cite{rawls1971,parfit1997,frankfurt1987,sen1979}. If a planner has to encode
\emph{some} distributional rule, the interesting engineering question is what happens
when we encode the alternatives explicitly and compare them.

\subsection{Contribution}

This paper takes five established principles of distributive justice --- Equal,
Utilitarian, Sufficientarian, Prioritarian and Proportional --- from their statements in
the philosophical literature, encodes each as an explicit constraint system over a
discrete multi-aircraft routing problem, and solves the result with an optimising SMT
solver~\cite{demoura2008z3,bjorner2015nuz}. The five are then compared on the routings
they produce, the noise outcomes they deliver, and what they cost to compute. The system
is not a thought experiment: it is the solver behind a participatory workshop tool in
which residents of Austin, TX select a principle and see the resulting flight paths over
a $5\times5$ map of their own city.

Our specific contributions are:

\begin{enumerate}
\item \textbf{An explicit encoding of five fairness principles} as hard constraints and
      optimisation objectives over a common routing model
      (Section~\ref{sec:encoding}), including a construction that makes ``benefit'' ---
      required by Proportional but absent from any noise model --- an editable input
      rather than a hidden constant.

\item \textbf{Evidence that the translation from words to constraints is not neutral},
      in three measured forms. Encoding Equal as a satisfaction constraint rather than
      the optimisation its statement also admits leaves it subsumed by Sufficientarian,
      so under the conditions we identify the difference a user sees is an arbitrary
      tie-break (Section~\ref{sec:mismatch}). Because noise adds in energy, expressing
      fairness in \emph{added} decibels makes already-loud communities some $70\times$
      cheaper to overfly (Section~\ref{sec:concavity}). And an ablation shows three of
      the five principles bind on nothing at the deployed setting
      (Section~\ref{sec:degeneracy}).

\item \textbf{A tractability comparison across principles.} The five differ by more than
      an order of magnitude in solve time, and not in the intuitive direction: the
      principles expressible as hard caps scale to fleets that leave the two
      decibel-valued objectives unable to prove optimality within a \SI{600}{\second}
      budget (Section~\ref{sec:tractability}).
\end{enumerate}

Throughout we report the awkward results alongside the clean ones. A tool that claims to
implement a fairness principle and in fact implements a near neighbour of it is worse
than one making no such claim; what follows is in part an audit of our own system.

\subsection{Related work}

The closest prior work is Gao et al.~\cite{gao2024noiseaware}, who formulate noise-aware
and equity-aware urban air traffic management as a multi-objective optimisation, trade
efficiency against fairness explicitly, and identify ambient-noise masking as a
mitigation lever --- the same asymmetry our Section~\ref{sec:concavity} measures. They
optimise a parameterised fairness criterion; we instead fix five named principles from
the philosophical literature, encode each separately, and ask which of them the solver
can tell apart. The two questions are complementary: theirs is how far efficiency and
equity trade, ours is whether a stated principle survives being written down.

On the normative side, Hauptvogel et al.~\cite{hauptvogel2021fairness} review aircraft
noise through social-justice theory and already enumerate distributive rules of the kind
Section~\ref{sec:principles} states. Their analysis stops at the enumeration; we take the
next step of committing each rule to a constraint system and measuring what it then does.

Compilation-based approaches to multi-agent path finding encode agent motion, collision
constraints and objectives as logical formulas discharged by a general
solver~\cite{surynek2019smtmapf,stern2019mapf}. We adopt that machinery but change the
coupling: our aircraft do not conflict with one another, and what links them is a shared,
nonlinear, spatially distributed \emph{externality}. Fair division studies allocation
under formal fairness criteria~\cite{moulin2003fair,walsh2020fairdivision} but typically
assumes the good is divisible or freely assignable; here the allocation must be
realisable as physical trajectories on a graph, which is what makes the encoding rather
than the axiomatics the hard part. Work on UAM noise itself has concentrated on source
characterisation, propagation and psychoacoustic
annoyance~\cite{nasaUAMnoise2020,kim2022indoor}; we take a coarse but standard acoustic
model as given.
